\begin{proposition}\label{prop:decic-no70}
An elliptic curve over a number field of degree $10$ has no point of
order $70$.
\end{proposition}
\begin{proof}
Put $X=X_1(70)$. The curve has genus $97$, good reduction at $3$,
and a rank-zero rational Jacobian by
\cite[Theorem~3.1(2)]{DEHMZ}. Let $D$ be the rational effective trace
divisor of a hypothetical level structure over a degree-$10$ field.
It has degree $10$, with repetitions allowed, and entirely
noncuspidal support.

First suppose that $\overline D$ is cuspidal. Lift this divisor in
the smooth cuspidal boundary over $\mathbb Q_3$, and call the lift
$C$. The class $[D-C]$ is torsion by rank zero and Manin--Drinfeld,
and reduces to zero. Torsion injectivity over the unramified odd
local field $\mathbb Q_3$ gives $D\sim C$ over $\mathbb Q_3$.
Since $D\ne C$, we have $h^0(D)\geq2$. Cohomology and base change
give the same inequality over $\mathbb Q$. This contradicts
Abramovich's bound~\cite{Abramovich}, since
\[
 \operatorname{gon}_{\mathbb C}X
 \geq\frac7{800}[\operatorname{PSL}_2(\mathbb Z):\overline{\Gamma_1(70)}]
 =\frac7{800}\,1728=\frac{378}{25}>10.
\]

The complete source census at $3$ gives
\[
\begin{array}{c|rrrrrrrrrr}
 \text{degree}&1&2&3&4&5&6&7&8&9&10\\ \hline
 \text{noncuspidal closed places}&0&0&0&0&0&48&408&735&2144&5880.
\end{array}
\]
These counts agree independently with Newton identities for the
degree-$194$ source Frobenius polynomial. The noncuspidal part of
$\overline D$ is therefore one closed point $P$ of degree
$d\in\{6,7,8,9,10\}$, with multiplicity one. The global and finite
cusp signatures at $3$ are both
\[
                         1^{24}4^6 6^4 12^2.
\]
Thus its cuspidal complement lifts to a rational cuspidal divisor
$C$ of degree $10-d$. In this step we use the equality of the
individual global and finite cusp orbits, checked by primitive
Tate-column arithmetic.

Let $Z=X_1(35)$, and consider the two rational maps
\[
 \begin{split}
 \pi_1(E,Q)&=(E,2Q),\\
 \pi_2(E,Q)&=(E/\langle35Q\rangle,\phi(2Q)).
 \end{split}
\]
The retained point in the second map has order $35$, since its
subgroup has odd order and the isogeny has degree $2$.
The Jacobian $J_Z(\mathbb Q)$ has rank zero, and its rational
torsion is generated by classes of rational cuspidal divisors
\cite[Theorems~3.1(2), 4.13 and Table~2]{DEHMZ}.

Work in the point-level convention and put
\[
                         A_2=T_2-1-2\langle2\rangle.
\]
This operator kills $J_Z(\mathbb Q)$ without any injectivity
assumption at $2$. Indeed, on geometric cusps the point-level
boundary formula is
\[
             T_2=\sigma_2+2\langle2\rangle\sigma_2^{-1},
\]
where $\sigma_2$ is cyclotomic Galois. A rational cuspidal divisor
is fixed by $\sigma_2$, so $A_2$ kills it as a divisor. Cuspidal
generation gives the assertion on $J_Z(\mathbb Q)$.

For a rational Tate cusp $c_0$, we also have $A_2c_0=0$ literally.
Applying $A_2$ to $[\pi_{i*}D-10c_0]$ shows that
$A_2\pi_{i*}D$ is principal. The divisor $\pi_{i*}C$ is rational
and cuspidal, so its $A_2$-image is zero as a divisor. Consequently,
for both $i=1,2$, the finite divisor
\[
                         A_2\pi_{i*}P                  \tag{70.1}
\]
must be principal.

We test (70.1) on the smooth genus-$25$ model of
$Z_{\mathbb F_3}$. The projective source diamond group has order
$12$. Complete normalization under simultaneous diamonds and
Frobenius gives respectively
\[
                         4,34,65,180,490
\]
source orbits in degrees $6,7,8,9,10$. Nontrivial stabilizers are
retained individually; no freeness assumption is made. For $770$
of the resulting $773$ representatives, the first map gives
$h^0(A_2\pi_{1*}P)=0$. The remaining three representatives, all of
degree $8$, give $h^0(A_2\pi_{2*}P)=0$. This contradicts (70.1).

Every Hecke divisor is computed from the three actual cyclic
$2$-kernels. A kernel orbit of length $e$ with target degree $f$
contributes the exact coefficient $de/f$. Projection norms and
coincident closed targets are included before signed divisors are
formed. In particular one source projection has norm $3$, and one
target coefficient is $3$; these are not replaced by unit weights.

The complete evidence is retained in \path{A2_all.json} and the five
original native RR packets. A fresh native replay reconstructs all
$3489$ prime references and all $776$ tested projected divisors,
including the three second projections. Independently, Sage verifies
each physical source level, its complete diamond orbit, all $2328$
geometric cyclic $2$-kernels, and the explicit V\'{e}lu coordinates.
All source, diamond and Hecke target keys agree with the independent
native prime keys after coincident targets are aggregated, with
every residue-degree norm unchanged. These checks are recorded in
\path{independent_genuine_isogenies.json} and the four files
\path{parent_native_s0.m}, \path{parent_native_s1.m},
\path{parent_native_s2.m}, \path{parent_native_s3.m}.
\end{proof}
